A la Doctora Zuleika Picó, médico radiólogo en Venezuela, por haber sido la inspiración
original de este Trabajo Fin de Máster y por dedicar su tiempo a validar,
desde su experiencia clínica diaria, la utilidad real del software aquí desarrollado.

A mi familia, por su constante acompañamiento, paciencia y apoyo emocional durante todo
el desarrollo de este trabajo.

A los profesores de la asignatura de Tecnologías del Habla y del Lenguaje Natural, 
Ramón López-Cózar, Zoraida Callejas y David Griol, por sus clases que me 
permitieron profundizar en la materia sobre la que gira este proyecto. Gracias a ellos
aprendí a la estructura de un sistema de procesamiento del lenguaje natural,
desde el reconocimiento de voz hasta la extracción de entidades, así como a utilizar
los Grandes Modelos de Lenguaje (LLMs) tanto para el desarrollo como para la evaluación
de la aplicación.

A mi tutor, Francisco Javier Melero Rus, por su guía fundamental en la definición, 
estructuración y alcance de este proyecto, y por su valiosa orientación sobre el desarrollo
de software en España, atendiendo rigurosamente a la normativa y a la legislación vigente.

Y al profesorado del máster en su conjunto, por proporcionarme las herramientas académicas 
y técnicas que me han permitido complementar con éxito el desarrollo de este Trabajo Fin 
de Máster.